\begin{table}[htbp]
\centering
\small
\caption{Sample construction and descriptive statistics.}
\label{tab:desc}
\begin{threeparttable}
\begin{tabular}{lcccc}
\toprule
 & Full sample & SD & Pre-shock & Post-shock \\
\midrule
\multicolumn{5}{l}{Panel A: Sample construction} \\
Firms in the 11 digital industries (2021Q1) & 384 & & & \\
\quad Less: no usable 2019 annual report & 86 & & & \\
\quad Final sample & 298 & & & \\
Firm-quarter observations & 6{,}900 & & 3,933 & 2,967 \\
High-TDDI firms & 100 & & & \\
\midrule
\multicolumn{5}{l}{Panel B: Main variables} \\
Log Tobin's Q & 0.861 & 0.558 & 0.901 & 0.809 \\
Tobin's Q & 2.804 & 1.899 & 2.960 & 2.597 \\
ROA & 0.003 & 0.037 & 0.006 & -0.000 \\
Gross margin & 0.360 & 0.231 & 0.372 & 0.344 \\
Selling-expense ratio & 0.117 & 0.121 & 0.115 & 0.120 \\
Period-expense ratio & 0.231 & 0.172 & 0.228 & 0.235 \\
Revenue growth & 0.123 & 0.534 & 0.174 & 0.073 \\
Firm size & 22.028 & 1.033 & 21.996 & 22.071 \\
Leverage & 0.361 & 0.192 & 0.347 & 0.378 \\
Cash ratio & 0.212 & 0.146 & 0.210 & 0.214 \\
TDDI & 0.001 & 0.998 & -0.001 & 0.004 \\
FPDCI & 0.008 & 1.003 & 0.011 & 0.004 \\
\bottomrule
\end{tabular}
\begin{tablenotes}[flushleft]
\footnotesize
\item Notes: All ratio variables are winsorized at the 1st and 99th percentiles. Pre-shock is 2018Q1--2021Q2; post-shock is 2021Q3--2023Q4. TDDI and FPDCI are standardized to mean zero and unit standard deviation. Firm size is the natural logarithm of total assets.
\end{tablenotes}
\end{threeparttable}
\end{table}
